\section{Основы C++}

\subsection{Каково назначение препроцессора языка C(++)? С какими сущностями он манипулирует? Имеет ли для него принципиальное значение синтаксис и семантика языка C++?}
\begin{definition}[Препроцессор]
\textit{Preprocessor} — это этап обработки исходного текста \emph{до} компиляции. Он выполняет текстовые преобразования: подключает файлы, разворачивает макросы, включает или исключает фрагменты кода по условиям.
\end{definition}

Назначение препроцессора:
\begin{itemize}
    \item подключение заголовков через \texttt{\#include};
    \item условная компиляция через \texttt{\#if}, \texttt{\#ifdef}, \texttt{\#ifndef} и т.д.;
    \item определение и подстановка макросов через \texttt{\#define};
    \item организация конфигурации кода под разные платформы, режимы сборки и окружения.
\end{itemize}

Препроцессор работает не с объектами языка C++ как таковыми, а с \textbf{текстом исходника}. Поэтому для него принципиально важны строки, токены и директивы, но \emph{не} семантика C++ в полном смысле: он не понимает типы, классы, перегрузку функций и другие правила компилятора.

\begin{remark}
Препроцессор можно рассматривать как текстовый механизм, а не как часть собственно языка C++. Именно поэтому ошибки препроцессора часто выглядят грубо и локально, а не как «умные» сообщения о типах.
\end{remark}

\subsection{Как на примере директивы \texttt{assert} препроцессор меняет код программы для конфигураций debug и release?}
Директива \texttt{\#include <cassert>} подключает проверку \texttt{assert}. Её поведение зависит от того, определён ли макрос \texttt{NDEBUG}.

\begin{definition}[Debug и release]
\begin{itemize}
    \item \textbf{debug} — режим отладки: проверки обычно включены;
    \item \textbf{release} — режим поставки: проверки обычно выключены ради скорости.
\end{itemize}
\end{definition}

\begin{example}
\begin{verbatim}
#include <cassert>

int f(int x) {
    assert(x != 0);
    return 100 / x;
}
\end{verbatim}
\end{example}

Если \texttt{NDEBUG} \emph{не} определён, то \texttt{assert(x != 0)} работает как runtime-check: при ложном условии программа аварийно завершится, обычно с диагностикой.  
Если \texttt{NDEBUG} определён, то макрос \texttt{assert} разворачивается в пустую конструкцию вида \texttt{((void)0)} и не влияет на runtime.

\begin{remark}
Это хороший пример принципа \textit{zero overhead}: проверки удобны в debug, но не должны накладывать стоимость в release.
\end{remark}

\subsection{Что содержит и что не содержит листинг (результат работы препроцессора)?}
\begin{definition}[Листинг]
\textit{Listing} — это текст, полученный после работы препроцессора, то есть уже «развёрнутый» исходник.
\end{definition}

Обычно листинг содержит:
\begin{itemize}
    \item подставленные заголовочные файлы;
    \item раскрытые макросы;
    \item результат условной компиляции;
    \item служебные директивы вида \texttt{\#line} или аналогичные, помогающие компилятору сопоставлять строки с исходными файлами.
\end{itemize}

Обычно листинг \emph{не содержит}:
\begin{itemize}
    \item самих директив \texttt{\#define}, \texttt{\#include} в их исходном виде;
    \item комментариев как самостоятельного материала;
    \item логики компиляции по типам и семантике C++.
\end{itemize}

\begin{remark}
Если посмотреть на листинг, можно увидеть, что препроцессор делает не «умное преобразование программы», а именно текстовую подстановку и фильтрацию фрагментов.
\end{remark}

\subsection{Что делают директивы условного включения \texttt{\#if}, \texttt{\#ifdef}, \texttt{\#ifndef}, \texttt{\#else}, \texttt{\#elif}, \texttt{\#endif}?}
Эти директивы позволяют включать или исключать фрагменты исходного текста в зависимости от условий препроцессора.

\begin{definition}[Условная компиляция]
\textit{Conditional compilation} — это механизм, при котором разные части кода присутствуют или отсутствуют в итоговом translation unit в зависимости от макросов и выражений препроцессора.
\end{definition}

Смысл директив:
\begin{itemize}
    \item \texttt{\#if <expr>} — включает блок, если выражение истинно;
    \item \texttt{\#ifdef X} — эквивалентно проверке «определён ли макрос \texttt{X}»;
    \item \texttt{\#ifndef X} — эквивалентно проверке «не определён ли \texttt{X}»;
    \item \texttt{\#else} — альтернативная ветка;
    \item \texttt{\#elif} — цепочка альтернативных условий;
    \item \texttt{\#endif} — завершение условного блока.
\end{itemize}

\begin{example}
\begin{verbatim}
#ifdef _WIN32
    // Windows-specific code
#elif defined(__linux__)
    // Linux-specific code
#else
    // Other platform
#endif
\end{verbatim}
\end{example}

\subsection{Как можно конфигурировать текст исходного кода программы условными директивами?}
Условные директивы применяют, когда один и тот же проект должен по-разному собираться на разных платформах или в разных режимах.

\begin{example}
\begin{verbatim}
#ifdef DEBUG
    std::cerr << "Debug mode\n";
#endif

#ifdef _WIN32
    // Windows implementation
#else
    // POSIX implementation
#endif
\end{verbatim}
\end{example}

Типичные случаи:
\begin{itemize}
    \item платформенно-зависимый код (\texttt{Windows}, \texttt{Linux}, \texttt{macOS});
    \item включение отладочных проверок;
    \item включение экспериментальных возможностей;
    \item защита от неподдерживаемой конфигурации.
\end{itemize}

\subsection{Для чего используется директива \texttt{\#include}? Чем отличаются системные и пользовательские заголовки? Каков порядок включения стандартных и пользовательских заголовков?}
Директива \texttt{\#include} вставляет содержимое другого файла в текущий translation unit на этапе препроцессинга.

\begin{definition}[Системный и пользовательский include]
\begin{itemize}
    \item \textbf{Системный (библиотечный) заголовок}: подключается как \texttt{\#include <...>};
    \item \textbf{Пользовательский (проектный) заголовок}: подключается как \texttt{\#include "..."};
\end{itemize}
\end{definition}

Разница в синтаксисе отражает намерение:
\begin{itemize}
    \item \texttt{<...>} — искать в системных путях и путях, заданных сборкой;
    \item \texttt{"..."} — сначала искать рядом с текущим файлом, затем в системных путях.
\end{itemize}

Стилевое соглашение обычно такое:
\begin{enumerate}
    \item сначала стандартные заголовки;
    \item затем пустая строка;
    \item потом проектные заголовки.
\end{enumerate}

\begin{remark}
Такой порядок делает зависимости более читаемыми и помогает быстро увидеть, где заканчивается стандартная библиотека и начинается код проекта.
\end{remark}

\subsection{В чём проблема повторного косвенного включения заголовочного файла? Привести пример.}
\begin{definition}[Повторное косвенное включение]
Повторное косвенное включение возникает, когда один и тот же заголовок попадает в translation unit несколько раз не напрямую, а через цепочку других заголовков.
\end{definition}

\begin{example}
Пусть:
\begin{verbatim}
// a.h
#include "common.h"

// b.h
#include "common.h"

// main.cpp
#include "a.h"
#include "b.h"
\end{verbatim}

Тогда \texttt{common.h} будет включён дважды, если не защитить его от повторного включения.
\end{example}

Проблема заключается в том, что без специальных мер повторное включение может привести к:
\begin{itemize}
    \item множественным определениям;
    \item повторному объявлению типов и функций;
    \item конфликтам имён;
    \item ошибкам компиляции.
\end{itemize}

\subsection{Что такое header guard и какую проблему он предотвращает? Привести пример.}
\begin{definition}[Header guard]
\textit{Header guard} — это классическая защита заголовочного файла от повторного включения.
\end{definition}

Классическая реализация:
\begin{verbatim}
#ifndef PROJECT_FOO_BAR_H
#define PROJECT_FOO_BAR_H

// contents of foo/bar.h

#endif // PROJECT_FOO_BAR_H
\end{verbatim}

Смысл: при первом включении макрос ещё не определён, поэтому содержимое заголовка попадает в translation unit, а затем макрос определяется. При повторном включении содержимое пропускается.

\begin{example}
\begin{verbatim}
#ifndef MYPROJECT_UTILS_VECTOR_H
#define MYPROJECT_UTILS_VECTOR_H

struct Vec {
    double x, y;
};

#endif
\end{verbatim}
\end{example}

\subsection{Можно ли использовать заголовочный файл без header guard? Какие ограничения это накладывает?}
Технически можно, но только если вы \textbf{гарантируете}, что такой заголовок попадёт в каждый translation unit не более одного раза.

На практике это ограничение слишком жёсткое: при росте проекта один и тот же заголовок почти неизбежно начинает подключаться косвенно через разные файлы.

Если guard-ов (или \texttt{\#pragma once}) нет, то возникают риски:
\begin{itemize}
    \item повторные объявления и определения;
    \item ошибки компиляции из-за конфликтов;
    \item хрупкая зависимость от порядка \texttt{\#include};
    \item сложная поддержка и рефакторинг.
\end{itemize}

\begin{remark}
Поэтому в обычном C++-коде заголовок без защиты от повторного включения считают плохой практикой. Исключения бывают только в узких технических сценариях, где повторное включение задумано специально.
\end{remark}

\subsection{Как именуют header guard с привязкой к пути к файлу и имени проекта? Можно ли использовать \texttt{\#pragma once}?}
Стилевое соглашение для header guard обычно такое:
\begin{itemize}
    \item имя в \texttt{UPPER\_CASE};
    \item в него включают имя проекта;
    \item затем путь к файлу, нормализованный к символам \texttt{\_};
    \item часто добавляют суффикс \texttt{\_H} или \texttt{\_HPP}.
\end{itemize}

\begin{example}
Если файл лежит как \texttt{project/include/math/vector.h}, то guard может выглядеть так:
\begin{verbatim}
PROJECT_INCLUDE_MATH_VECTOR_H
\end{verbatim}
\end{example}

Вместо классического guard иногда используют:
\begin{verbatim}
#pragma once
\end{verbatim}

Это короткий и удобный способ запретить повторное включение. На практике он широко поддерживается компиляторами, но классический header guard остаётся максимально переносимым и стандартизованным способом.

\begin{remark}
Если в курсе требуется именно стиль с guard-ами, лучше уметь писать и понимать оба варианта.
\end{remark}

\subsection{Как включить в модуль системный заголовочный файл? Как включить заголовочный файл из проекта?}
В модульной модели C++ есть несколько вариантов подключения.

Если речь идёт о \textbf{module interface unit} или \textbf{module implementation unit}, то стандартные заголовки обычно подключают либо через \texttt{\#include}, либо через импорт стандартной библиотеки, если это поддерживается инструментами и стандартной библиотекой в данной среде.

\begin{example}
\begin{verbatim}
module;

#include <vector>
#include "project/config.h"

export module foo;

export int f();
\end{verbatim}
\end{example}

Здесь:
\begin{itemize}
    \item \texttt{\#include <vector>} — системный заголовок;
    \item \texttt{\#include "project/config.h"} — заголовок проекта;
    \item строка \texttt{module;} означает \textit{global module fragment}, где разрешены обычные include перед объявлением модуля.
\end{itemize}

\subsection{Что обычно кладут в заголовочный файл — интерфейсную часть некоторого модуля \texttt{foo}? А что обычно не кладут — и почему?}
В заголовок обычно кладут только \textbf{интерфейс}:
\begin{itemize}
    \item объявления функций;
    \item определения классов и структур;
    \item объявления шаблонов;
    \item константные объявления и alias-ы;
    \item public API, которым пользуется внешний код.
\end{itemize}

Обычно не кладут:
\begin{itemize}
    \item тяжёлую реализацию алгоритмов;
    \item внутренние helper-функции;
    \item подробности хранения данных;
    \item всё, что можно скрыть, не ломая API.
\end{itemize}

Причины:
\begin{enumerate}
    \item уменьшить связанность кода;
    \item сократить время компиляции;
    \item скрыть детали реализации;
    \item избежать повторного включения и проблем с множественными определениями.
\end{enumerate}

\begin{remark}
Хороший заголовок отвечает на вопрос «\emph{что умеет модуль}», а не «\emph{как он это делает}».
\end{remark}

\subsection{Чем в C++ описание чего-либо принципиально отличается от определения? Как наблюдать эффект описания и эффект определения?}
\begin{definition}[Declaration]
\textit{Declaration} — это описание сущности для компилятора: мы сообщаем, что такая сущность существует и как с ней можно работать.
\end{definition}

\begin{definition}[Definition]
\textit{Definition} — это определение сущности, то есть создание конкретного объекта, функции или другого элемента программы.
\end{definition}

Формулировка в простом виде:
\begin{itemize}
    \item \textbf{описать} — значит сообщить компилятору о существовании сущности;
    \item \textbf{определить} — значит реально создать сущность, то есть выделить для неё память или дать тело функции.
\end{itemize}

\begin{example}
\begin{verbatim}
extern int x;   // declaration
int x = 5;      // definition
\end{verbatim}
\end{example}

Здесь первая строка сообщает о переменной \texttt{x}, но не создаёт её. Вторая строка уже создаёт объект \texttt{x} в памяти.

Для функций:
\begin{verbatim}
int sum(int, int);      // declaration
int sum(int a, int b) { // definition
    return a + b;
}
\end{verbatim}

\begin{remark}
Эффект описания наблюдается в том, что компилятор начинает «знать имя» и тип сущности. Эффект определения наблюдается в том, что сущность получает реальное существование в программе: переменная начинает занимать память, а функция — иметь тело, которое можно вызвать и слинковать.
\end{remark}
